% main.tex

\documentclass[12pt]{article}

% preamble.tex

\usepackage[utf8]{inputenc}
\usepackage[T1]{fontenc}
\usepackage[brazil]{babel}

\usepackage{amsmath, amssymb, amsfonts}
\usepackage{mathtools}

\usepackage{amsthm}
\theoremstyle{definition}
\newtheorem{problemstatement}{Enunciado do Problema}

\usepackage{geometry}
\geometry{a4paper, margin=1in}

\usepackage{listings}
\usepackage{xcolor}

% Cores personalizadas para o Código Fonte (listings)
\definecolor{codegreen}{rgb}{0,0.6,0}
\definecolor{codegray}{rgb}{0.5,0.5,0.5}
\definecolor{codepurple}{rgb}{0.58,0,0.82}
\definecolor{backcolour}{rgb}{0.95,0.95,0.95}

% Configuração do listings (Python)
\lstdefinestyle{mystyle}{
    language=Python,
    backgroundcolor=\color{backcolour},
    commentstyle=\color{codegreen},
    keywordstyle=\color{blue},
    numberstyle=\tiny\color{codegray},
    stringstyle=\color{codepurple},
    basicstyle=\ttfamily\footnotesize,
    breakatwhitespace=false,
    breaklines=true,
    captionpos=b,
    keepspaces=true,
    numbers=left,
    numbersep=5pt,
    showspaces=false,
    showstringspaces=false,
    showtabs=false,
    tabsize=2,
    frame=single,
    rulecolor=\color{black},
    % escapeinside={\%*}{*)}
}
\lstset{style=mystyle}

\usepackage{graphicx}

\usepackage{hyperref}
\hypersetup{
    colorlinks=true,
    linkcolor=blue!70!black,
    citecolor=green!60!black,
    urlcolor=red!70!black
}


\title{Problemas Matemáticos do Project Euler para Estudo}
\author{Fabio Siqueira}
\date{3 de junho de 2025 -- \today}

\begin{document}

% --- CAPA ---
\begin{titlepage}
    \centering
    \vspace*{3cm}
    
    % Título
    {\Huge \bfseries \textsc{Problemas Matemáticos do Project Euler para Estudo}\par}
    \vspace{0.6cm}

    % Subtítulo
    {\Large \textit{Registro Pessoal de Problemas Matemáticos e Aprendizados}\par}
    \vspace{3cm}

    % Autor
    {\LARGE Fabio Siqueira\par}
    \vspace{1.5cm}

    % Data
    {\large 3 de junho de 2025 -- \today\par}

    \vfill

    {\small \textit{"Ad scientiae complementum oportet omnia et singula, quae ad institutum nostrum pertinent, continuo et nullibi interrupto cogitationis motu perlustrare, atque illa sufficienti et ordinata enumeratione complecti."} \par
    \vspace{0.1cm}
    -- René Descartes, \textit{Regulae ad directionem ingenii} (Regula VII)\par}

    \vspace{2cm}
\end{titlepage}



% --- INTRODUÇÃO ---
\section{Introdução}
\label{sec:introducao}

Este livro é um projeto que nasceu da minha curiosidade e interesse pela matemática.

Atualmente, a principal área de pesquisa na minha vida é a da tecnologia.
Portanto, este projeto é o meu primeiro passo de estudo autodidata sobre matemática pura na vida.

Nele, você vai me encontrar resolvendo problemas que encontrei no website \textsc{Project Euler} (\url{https://projecteuler.net/}).

\newpage



% --- SUMÁRIO ---
\tableofcontents

\newpage



% --- PROBLEMAS ---
\section{Multiples of 3 or 5 (Múltiplos de 3 ou 5)}
\label{sec:problem1}
% --- problems/problem0001.tex ---

A resolução do problema consiste em encontrar a soma de todos os números naturais abaixo de 1000 que são múltiplos de 3 ou 5.

\subsection*{Enunciado Oficial (Inglês)}
\begin{quote}
If we list all natural numbers below 10 that are multiples of 3 or 5, we get 3, 5, 6, and 9. The sum of these multiples is 23.
Find the sum of all multiples of 3 or 5 below 1000.
\end{quote}

\subsection*{Enunciado do Problema (Português)}
\begin{quote}
Se listarmos todos os números naturais abaixo de 10 que são múltiplos de 3 ou 5, obtemos 3, 5, 6 e 9. A soma desses múltiplos é 23.
Encontre a soma de todos os múltiplos de 3 ou 5 menores que 1000.
\end{quote}



\subsection*{Abordagem Matemática}

Para encontrar a soma de todos os números naturais abaixo de 1000 que são múltiplos de 3 ou 5, precisamos de um método que evite contar os números duas vezes.

O Princípio da Inclusão-Exclusão é perfeitamente adequado para esta tarefa.
Ele dita que somamos os múltiplos de 3, adicionamos a soma dos múltiplos de 5 e, em seguida, subtraímos a soma dos números que são múltiplos de ambos (ou seja, múltiplos de $3 \times 5 = 15$).

Isso pode ser expresso de forma concisa como:
\[
\sum_3 + \sum_5 - \sum_{15}
\]
Neste contexto, $\sum_3$, $\sum_5$ e $\sum_{15}$ representam, respectivamente, a soma de todos os múltiplos de 3, de 5 e de 15, que são estritamente menores que 1000.

Para calcular cada uma dessas somas individuais ($\sum_x$), usamos a conhecida fórmula para a soma de uma progressão aritmética.
Especificamente, a soma dos múltiplos de um número $x$ até um certo limite $n$ é dada por:
\[
x \cdot \frac{m(m + 1)}{2}, \quad \text{onde } m = \left\lfloor \frac{n}{x} \right\rfloor
\]
Aqui, $m$ é a contagem de quantos múltiplos de $x$ existem até $n$.

Para este problema, estamos interessados apenas nos números que estão abaixo de 1000. Isso significa que o maior número que consideramos em nossas somas é $n=999$. Aplicando isso à nossa fórmula para cada componente:
\[
\text{Soma} = 3 \cdot \frac{333 \cdot 334}{2} + 5 \cdot \frac{199 \cdot 200}{2} - 15 \cdot \frac{66 \cdot 67}{2} = \boxed{233168}
\]
Este cálculo resulta na soma total de todos os números naturais únicos abaixo de 1000 que são múltiplos de 3 ou 5.



\subsection*{Código Python}
O código Python a seguir itera pelos números abaixo de 1000 e soma aqueles divisíveis por 3 ou 5.
\begin{lstlisting}[caption=Solução do Problema 1 em Python, language=Python]
limit = 1000
total_sum = 0
for n in range(1, limit):
....if n % 3 == 0 or n % 5 == 0:
........total_sum += n
# print(total_sum)
\end{lstlisting}



\subsection*{Resposta}
A soma de todos os múltiplos de 3 ou 5 abaixo de 1000 é $\mathbf{233168}$.

\newpage


\section{Even Fibonacci Numbers (Números Pares de Fibonacci)}
\label{sec:problem2}
% --- problems/problem0002.tex ---

A resolução do problema consiste em encontrar o resultado da soma de todos os números pares que fazem parte da sequência Fibonacci e que não excedem o número 4.000.000 (quatro milhões).

\subsection*{Enunciado Oficial (Inglês)}
\begin{quote}
Each new term in the Fibonacci sequence is generated by adding the previous two terms. By starting with 1 and 2, the first 10 terms will be:

$1, 2, 3, 5, 8, 13, 21, 34, 55, 89, ...$

By considering the terms in the Fibonacci sequence whose values do not exceed four million, find the sum of the even-valued terms.
\end{quote}

\subsection*{Enunciado do Problema (Português)}
\begin{quote}
Cada novo termo na sequência de Fibonacci é gerado pela adição dos dois termos anteriores. Começando com 1 e 2, os 10 primeiros termos serão:

$1, 2, 3, 5, 8, 13, 21, 34, 55, 89, ...$

Considerando os termos na sequência de Fibonacci cujos valores não excedem quatro milhões, encontre a soma dos termos de valor par.
\end{quote}

\subsection*{Abordagem Matemática}
Para resolver este problema, precisamos gerar os termos da sequência de Fibonacci até que eles ultrapassem o valor de quatro milhões. A sequência, conforme especificado, inicia com os termos 1 e 2.
\\A definição recursiva da sequência de Fibonacci é:
\[
\boxed{F_n = F_{n-1} + F_{n-2}}
\]
\\Para este problema, temos as condições iniciais:
\[
\boxed{F_1 = 1} \text{ e } \boxed{F_2 = 2}
\]
\\Os termos subsequentes são, por exemplo:
\begin{align*}
F_3 &= F_2 + F_1 = 2 + 1 = 3 \\
F_4 &= F_3 + F_2 = 3 + 2 = 5 \\
F_5 &= F_4 + F_3 = 5 + 3 = 8 \\
F_6 &= F_5 + F_4 = 8 + 5 = 13 \\
F_7 &= F_6 + F_5 = 13 + 8 = 21 \\
F_8 &= F_7 + F_6 = 21 + 13 = 34
\end{align*}
e assim por diante.

Enquanto geramos cada termo, verificamos duas condições:
\begin{enumerate}
    \item O valor do termo não deve exceder 4.000.000.
    \item Se o termo for par, ele deve ser adicionado a uma soma acumuladora.
\end{enumerate}

Eis o início do processo de cálculo da soma dos termos pares:
\begin{itemize}
    \item $F_1 = 1$ (ímpar)
    \item $F_2 = 2$ (par) Soma = $\mathbf{2}$
    \item $F_3 = 3$ (ímpar)
    \item $F_4 = 5$ (ímpar)
    \item $F_5 = 8$ (par) Soma = $2 + 8 = \mathbf{10}$
    \item $F_6 = 13$ (ímpar)
    \item $F_7 = 21$ (ímpar)
    \item $F_8 = 34$ (par) Soma = $10 + 34 = \mathbf{44}$
    \item $F_9 = 55$ (ímpar)
    \item $F_{10} = 89$ (ímpar)
    \item $F_{11} = 144$ (par) Soma = $44 + 144 = \mathbf{188}$
    \item \dots
\end{itemize}

Este mesmo procedimento deve ser repetido diversas vezes, adicionando cada termo par encontrado à uma soma parcial, até que o termo de Fibonacci gerado exceda quatro milhões, nos dando a soma final.

Apresentar a lista completa de todos os termos e a soma manual passo a passo até o final seria inviável por conta da quantidade de termos. É neste ponto que um algoritmo, como o implementado em Python, se torna essencial para calcular a soma final de forma eficiente. Obrigado, computadores!

Uma observação interessante é que na sequência de Fibonacci os termos pares \\($F_2, F_5, F_8, F_{11}, \dots$)\\ ocorrem em um padrão regular: a cada três termos. Embora esta propriedade possa ser explorada para otimizar o cálculo (gerando apenas os termos pares), a abordagem de gerar cada termo e testar sua paridade é direta e fácil de compreender.

\subsection*{Código Python}
O código Python abaixo implementa a geração da sequência de Fibonacci e a soma dos termos pares que atendem à condição do problema.
\begin{lstlisting}[caption=Solução do Problema 2 em Python, language=Python]
# Sequencia de Fibonacci comecando com 1 e 2
termo_a = 1
termo_b = 2

soma_pares = 0
limite = 4000000

# Iteramos enquanto o termo atual (termo_b) nao excede o limite
while termo_b <= limite:
    if termo_b % 2 == 0:  # Verifica se o termo atual e' par
        soma_pares += termo_b
    
    # Gera o proximo termo da sequencia
    proximo_termo = termo_a + termo_b
    termo_a = termo_b          # Atualiza o termo anterior
    termo_b = proximo_termo    # Atualiza o termo atual

# print(soma_pares) # O resultado esperado e' 4613732
\end{lstlisting}

\subsection*{Resposta}
A soma dos termos pares da sequência de Fibonacci, cujos valores não excedem quatro milhões, é $\mathbf{4613732}$.

\newpage


\section{Largest Prime Factor (Maior Fator Primo)}
\label{sec:problem3}
% --- problems/problem0003.tex ---

Problema ainda não resolvido.

\newpage


% \section{Problema X}
%\label{sec:problemX}
% \input{problems/problem000X.tex}

\end{document}
